Le but de cette partie est l'étude d'une transformation spontanée dans une
pile.

On considère la pile Zinc/Argent. Cette pile est constituée des éléments
suivants~:
\begin{itemize}
    \item Un bécher contenant une solution aqueuse de nitrate d'argent
          $\ce{Ag^{+}_{(aq)} + NO^{-}_{3(aq)}}$ de volume $V_{1}$ et de
          concentration molaire $C_{1}$~;
    \item Un bécher contenant une solution aqueuse de nitrate de zinc
          $\ce{Zn^{2+}_{(aq)} + 2NO^{-}_{3(aq)}}$de volume $V_{2}$ et de
          concentration molaire $C_{2}$~;
    \item Un fil d'argent $\ce{Ag_{(s)}}$~;
    \item Une plaque mince du zinc $\ce{Zn_{(s)}}$.
    \item Un pont salin.
\end{itemize}

\begin{donnees}
    \begin{table}[H]
    \centering
    \begin{tabularx}{\textwidth}{|C|C|C|}
        \hline
        $C_{1}=\qty{2.0e-1}{\mol\per\L}$ & $C_{2}=\qty{2.0e-2}{\mol\per\L}$ &
        $1 \mathcal{F}=\qty{9.65e4}{\C\per\mol}$ \\
        \hline
        \multicolumn{3}{|c|}{La constante d'équilibre associée à l'équation~:
        $\ce{2Ag^{+}_{(aq)} + Zn_{(s)} <=>[(1)][(2)] 2Ag_{(s)} +
        Zn^{2+}_{(aq)}}$ est~: $\mathrm{K}=10^{52}$} \\
        \hline
        \end{tabularx}
    \end{table}
\end{donnees}

On branche, en série aux bornes de la pile, un ampèremètre et un conducteur
ohmique. Le circuit est alors traversé par un courant électrique.
\begin{enumerate}
    \item Déterminer la valeur du quotient de réaction $Q_{r,i}$, du système
          chimique à l'état initial .
    \item Déduire, en justifiant votre réponse, le sens d'évolution spontané du
          système chimique lors du fonctionnement de la pile.
    \item On laisse la pile fonctionner pendant une durée très longue jusqu'à ce
          qu'elle s'épuise.

          Déterminer la valeur de la quantité d'électricité maximale $Q_{\max}$,
          qui a traversé le conducteur ohmique du début de fonctionnement de la
          pile jusqu'à ce qu'elle s'épuise sachant que l'avancement maximale est
          $x_{\max}=\qty{5e-3}{\mol}$.
\end{enumerate}

